\section{Components}

\begin{frame}[fragile]{What Is an OWL Component?}
	\begin{itemize}
		\item A component is a class + template (+ optional props/state).
		\item Parents compose children to build complex screens.
\begin{block}{Shape}
\begin{lstlisting}[style=js]
export class StudentCard extends Component {
	static template = "school_manager.StudentCard";
	static props = ["student"];
	static components = {};
}
\end{lstlisting}
\end{block}
	\end{itemize}
\end{frame}

\begin{frame}[fragile]{State with \texttt{useState}}
\begin{lstlisting}[style=js]
import { Component, useState } from "@odoo/owl";

export class StudentFilter extends Component {
	static template = "school_manager.StudentFilter";
	setup() {
		this.state = useState({
			query: "",
			onlyActive: true,
		});
	}
	onQueryInput(ev) {
		this.state.query = ev.target.value;
	}
}
\end{lstlisting}
	\begin{itemize}
		\item Changing reactive state re-renders the component.
	\end{itemize}
\end{frame}

\begin{frame}[fragile]{Props and Child Components}
\begin{lstlisting}[style=js]
export class StudentList extends Component {
	static template = "school_manager.StudentList";
	static components = { StudentCard };
	static props = {
		students: { type: Array },
	};
}
\end{lstlisting}
\begin{lstlisting}[style=xml]
<t t-name="school_manager.StudentList">
  <t t-foreach="props.students" t-as="student" t-key="student.id">
	<StudentCard student="student"/>
  </t>
</t>
\end{lstlisting}
\end{frame}

\begin{frame}[fragile]{Events / Callbacks}
\begin{lstlisting}[style=js]
export class StudentCard extends Component {
	static props = {
		student: { type: Object },
		onSelect: { type: Function, optional: true },
	};
	select() {
		this.props.onSelect?.(this.props.student);
	}
}
\end{lstlisting}
	\begin{itemize}
		\item Prefer callbacks/props over tightly coupled parent lookups.
	\end{itemize}
\end{frame}

\begin{frame}{Component Best Practices}
	\begin{itemize}
		\item One component = one UI responsibility.
		\item Declare \texttt{static props} clearly.
		\item Use \texttt{t-key} in loops.
		\item Keep templates close to the component's module assets.
	\end{itemize}
\end{frame}
